\renewcommand{\theposttitle}{Creative Resistance and Finding Purpose}
\renewcommand{\thepostdate}{2022-11-01}
\vspace*{\fill}
\begin{center}
\includegraphics[width=\textwidth, height=0.55\textheight, keepaspectratio]{C:/Users/James/Documents/GitHub/VisualIdeasBook/data/images/11f5061c3d3677c03624e5319326bf9d.png}
\end{center}
\vspace{2mm}
\begin{center}
\begin{minipage}{0.92\textwidth}
\small
Where we place our focus affects the range and depth of our expertise.

As generalists develop range by spreading their focus over many skills, research suggests they are better creative, out-of-the-box thinking, because they see dots from other domains unseen by the specialist.

As specialists develop filling singular skillsets with depth, they are more efficient and can probably better innovate on their given area of expertise.

There're situations where generalists shine and situations where specialists shine.

Where do you stand on the generalist vs specialist argument?
\end{minipage}
\end{center}
\vspace*{\fill}